\section{Структура экзамена}

\subsection{Форма и регламент}
\begin{itemize}
    \item Экзамен проводится в письменной форме.
    \item Продолжительность: 1 час.
    \item Работа выполняется самостоятельно.
\end{itemize}

\subsection{Структура работы и критерии оценивания}
Максимальное количество баллов за работу --- \textbf{30}.

\begin{enumerate}
    \item \textbf{Теоретическая часть (25 баллов):}
    \begin{itemize}
        \item \textbf{2 определения.} Максимум --- \textbf{5 баллов}.
        \item \textbf{2 формулировки теорем.} Максимум --- \textbf{10 баллов}.
        \item \textbf{Доказательство одной теоремы.} Максимум --- \textbf{10 баллов}. \\
        \emph{Важно:} Студент самостоятельно выбирает, какую из попавшихся в билете теорем он будет доказывать.
    \end{itemize}

    \item \textbf{Практическая часть (5 баллов):}
    \begin{itemize}
        \item \textbf{2 задачи.} Максимум --- \textbf{5 баллов}. \\
        \emph{Темы задач:}
        \begin{enumerate}
            \item[-] Условия Коши—Римана.
            \item[-] Разложение функции в ряд (Тейлора или Лорана).
            \item[-] Вычисление интегралов от функции комплексного переменного (непосредственное, с помощью теоремы Коши или с помощью вычетов).
        \end{enumerate}
    \end{itemize}
\end{enumerate}

\subsection{Минимальный порог}
Для получения положительной оценки необходимо набрать \textit{не менее 10 баллов}.